\section{Discussion}
\subsection{What a Centering Gain Measures}
The additive model accounts for the observations across two datasets, two regimes, and eight transforms. Centering is an estimate of a participant offset, its benefit is set by the participant share of discriminant variance, and its cost is the noise of the estimate, $\Sigma_e/m$. Single trials sit where the share is small. There the full gain is between 0.000 and 0.036, a reference set of one or two trials lowers accuracy, and the break-even size of Corollary~\ref{cor:breakeven} locates the crossing in all five paradigms from two numbers, the uncentered and the fully centered accuracy. Trial-averaged ERPs sit where the share is large, because averaging divides $\Sigma_e$ by the number of trials and leaves $\Sigma_a$ unchanged. There a panel of two reference observations already carries most of the gain.

\subsection{Offsets and Gains}
Second-order alignment and first-order centering act on different parts of the participant effect. On single trials Euclidean alignment changes accuracy by 0.010 or less in every paradigm, with intervals that include zero, while centering gives a positive gain in four of five. On trial-averaged ERPs alignment alone recovers about 70\% of the centering gain, which the additive model does not predict. A participant gain, a multiplicative scale on the whole response, is therefore a substantial part of between-participant variability in averaged ERPs. Centering removes the offset and leaves a gain in place, alignment divides out the gain and leaves the offset, and on these data the offset is the larger part ($+0.066$ [$+0.036$, $+0.096$] for centering over alignment). The two families are complementary in the model, and on averaged data their combination adds no measurable accuracy over centering alone.

\subsection{Reporting Transductive Results}
A decoding accuracy obtained with any participant-level transform should state three things: which observations of the test participant the transform used, whether they include the observation being classified, and whether their class composition was fixed by the design. With a complete balanced panel, cross-participant classification becomes a within-participant contrast with separation $K/(K-1)$ times the asymptote, and the joint-assignment readout raises accuracy further by using the known composition. Such a result describes a repeated-measures protocol in which every test participant is recorded in every condition. It is a different quantity from the accuracy of a decoder applied to one new observation of a new participant, and the two belong in separate columns of any comparison.

\subsection{Scope and Future Directions}
The present work characterizes linear readouts of time-domain ERP features in six paradigms. The law is stated for a common class composition across participants, and Proposition~\ref{prop:law} predicts a participant-specific shift when the composition varies, which the oddball paradigms with resampled class proportions can measure directly. Extending the reference-size analysis to covariance-domain alignment, where the reference estimate is a matrix and its error grows with the channel count, would give Euclidean and Riemannian alignment their own break-even sizes. Motor-imagery and steady-state paradigms, with oscillatory features, are a natural next test of how the participant share $r$ varies across signal types.
